\begin{pastexam}[label={px:ncu-108-dsa-03}]{中央資工所 108 學年度 第 3 題}
（單選；每題 5 分，答錯倒扣 1 分）

Which of the following tree is a legal max-heap?
\begin{enumerate}[label=(\Alph*)]
\item A tree with level-order traversal sequence \{67, 45, 19, 22, 43, 16\}
\item A tree with level-order traversal sequence \{67, 43, 19, 22, 45, 16\}
\item A tree with level-order traversal sequence \{14, 18, 27, 19, 63, 48\}
\item A tree with level-order traversal sequence \{45, 22, 67, 8, 34, 52\}
\item None of the above.
\end{enumerate}
\end{pastexam}

\begin{pastsoln}
答案：(A)。

堆積是完整二元樹，層序就是它的陣列表示（資結書第~8.1~節）：第 $i$ 個的子節點是第 $2i$、$2i+1$ 個。
六個節點時，要檢查的父子關係是 $1$--$2$、$1$--$3$、$2$--$4$、$2$--$5$、$3$--$6$；最大堆積要求每個父節點都
不小於子節點。
\begin{itemize}
\item (A) $67 \ge 45, 19$；$45 \ge 22, 43$；$19 \ge 16$。 全部成立，是最大堆積。
\item (B) 第 $2$ 個 $43$ 的子節點有第 $5$ 個 $45 > 43$，不是。
\item (C) 根 $14$ 比子節點 $18, 27$ 都小，不是最大堆積。 （它其實是最小堆積：$14 \le 18, 27$；
$18 \le 19, 63$；$27 \le 48$。）
\item (D) 根 $45$ 的右子 $67 > 45$，不是。
\item (E) (A) 是，所以錯。
\end{itemize}
\end{pastsoln}
